\section{Numerical Simulations}\label{simulations}

Consider five planar second-order agents. The initial and target
configurations are respectively
$[y_1(0),\ldots,y_5(0)] = \left[\begin{smallmatrix} -3 & -3 & -3 & -3 & -3 \\ 0.002 & -1.15 & -0.55 & 0.55 & 1.15 \end{smallmatrix}\right]$ and $[y_1^\star,\ldots,y_5^\star] = \left[\begin{smallmatrix} 3 & 3 & 3 & 3 & 3 \\ 0 & -0.55 & -1.15 & 1.15 & 0.55 \end{smallmatrix}\right]$.
All agents start at rest, with $v_i(0)=\mathbf 0_n$ and
$\mathbf y^i(0)=y(0)$ for all $i\in\mathcal V$.
The cost of player $i$ is
$
J_i(y_i,y_{-i})
=
\frac12\|y_i-p_i\|^2
+
\frac{0.03}{2}
\sum_{j\ne i}\|y_i-y_j\|^2,
$
where
$
[p_1,\ldots,p_5]
=
[y_1^\star,\ldots,y_5^\star]B$ with $
B=1.15I_5-0.03\,\mathbf 1_5\mathbf 1_5^\top.
$
The pseudo-gradient is $1$-strongly monotone and
$1.15$-Lipschitz, yielding the unique unconstrained NE
$y^\star$. The extended pseudo-gradient satisfies
$\tilde l\simeq1.12161$.
The workspace contains three circular obstacles with centers
$o_1=(0.24,0)^\top$ and
$o_{2,3}=(1.72,\pm1.1)^\top$, and radii
$r_1=0.625$ and $r_2=r_3=0.315$.
We set the sensing radius to $d_c=2.2$ and the collision
distance to $d_{ij}=0.345$.
The collision and obstacle barriers are
$
	h_{ij}^{\mathrm{int}}
	=\|y_i-y_j\|^2-d_{ij}^2$,
	$h_{i,k}^{\mathrm{obs}}
	=\|y_i-o_k\|^2-r_k^2,
$
and the normalized matrix-tree certificate is
$
h^{\mathrm c}(y)
=
\frac{5}{512}\det L_\chi^{(r)}(y)-0.04$
with $
\underline{\lambda}_\chi=0.04,
$
where $L_\chi^{(r)}(y)$ is the principal cofactor
defined in Remark~\ref{tree}.

The ECBF gains are selected as
$
K_1=\alpha_1\alpha_2=4$ and $
K_2=\alpha_1+\alpha_2=4.6$,
with $\alpha_1\simeq1.16422$ and
$\alpha_2\simeq3.43578$.
The remaining algorithm parameters are
$
(c,\gamma,\xi,\epsilon,\nu)
=
(200,65,1.4,5\times10^{-6},0.02),
$
together with $\theta_i=50$ and $\delta_{p,i}=10^{-3}$.
The uniform tightening margins are
$
(\underline{\Delta}_{\mathrm{int}},
\underline{\Delta}_{\mathrm{obs}},
\Delta^{\mathrm c})
=
(0.5,0.7,0.4).
$
The physical ECBF terms and tightening margins
are scaled by $0.06$ for collision and obstacle
constraints and by $0.1$ for connectivity constraints,
while the mismatch terms remain unscaled.
For instance, the scaled local collision row is
$$
g_{ij,\mathrm{sc}}^i
=\chi_{ij}(y)\left[
0.06\left(q_{ij,i}^{\mathrm{int}}
+\frac12\Delta_{ij}^{\mathrm{int}}\right)
+\zeta_{ij}^i-\zeta_{ij}^j
\right]\leqslant0.
$$
The reference QP uses the same scaled constraints.
We take
$$
\mathcal D_R=
\left\{
(y,v)\ \middle|\
\begin{array}{l}
	y_i\in[-3.1,3.1]\times[-1.3,1.3],\\
	\max_{i<j}\|y_i-y_j\|\leqslant2.5,\\
	\max_i\|v_i\|\leqslant0.5
\end{array}
\right\}.
$$
The corresponding coefficient bounds are
$
(\bar a_{\mathrm{int}},
\bar a_{\mathrm{obs}},
\bar a_{\mathrm c})
=
(5,10.7689182372,2.69279673023).
$
The sensing-edge reactivation condition satisfies
$
2d_c\bar v_{ij,R}
\leqslant4.4
<
5.49625
=
\alpha_1(d_c^2-d_{ij}^2).
$
Both trajectories remain in this domain, with maximum
pairwise distance $2.3$ and maximum individual speed
$0.0874257$. Both runs are simulated 
using \texttt{ode15s}, with
$\mathrm{RelTol}=2\times10^{-9}$,
$\mathrm{AbsTol}=10^{-12}$, and
$\mathrm{MaxStep}=0.25$s.

We choose the design budgets
$r_u=0.061$ and $\delta_{\mathrm c}=0.03$, with
$
\mathcal U_R^u=\{u\mid\|u\|\leqslant15\}.
$
Along the tightened trajectory, define the trajectory-wise connectivity-estimation discrepancy as
$
\delta_{\mathrm c}^{\mathrm{tr}}(t)
=
\max_{u\in\mathcal U_R^u}
|
q^{\mathrm c}(\vartheta_{\mathrm{tight}}(t),u)
-
\sum_{i=1}^N
\hat q_i^{\mathrm c}(\vartheta_{\mathrm{tight}}(t),u_i)
|.
$
Table~\ref{tab:numerical-checks} summarizes
the numerical checks of the design and
initialization conditions.

\begin{table}[!t]
\caption{Selected Numerical Checks for
	Theorems~\ref{nominal_LES}--\ref{thm:closed_loop_LES}}
	\label{tab:numerical-checks}
	\centering
	\footnotesize
	\setlength{\tabcolsep}{2.5pt}
	\renewcommand{\arraystretch}{1.12}
	\begin{tabularx}{\columnwidth}{
			@{}
			>{\raggedright\arraybackslash}p{0.3\columnwidth}
			>{\centering\arraybackslash}p{0.3\columnwidth}
			>{\raggedright\arraybackslash}X
			@{}
		}
		\toprule
		Quantity & Value & Threshold\\
		\midrule
		
		\multicolumn{3}{@{}l}{
			\textit{Theorem \ref{nominal_LES}}
		}\\
		\midrule
		
		$c$
		& $200$
		& $>185.085$\\
		
		$\gamma$
		& $65$
		& $>56.6897$\\
		
		\shortstack[l]{
			$\min_{t}\lambda_2(L_\chi(y(t)))$
		}
	& $1.58616$
	& $>\underline{\lambda}_\chi=0.04$\\
		
		\addlinespace[2pt]
		\multicolumn{3}{@{}l}{
			\textit{Theorem \ref{thm:regional_safety}}
		}\\
		\midrule
		
		$\sqrt{2}\,\bar a_{\mathrm{int}}r_u$
		& $0.431335$
		& $<\underline{\Delta}_{\mathrm{int}}=0.5$\\
		
		$\bar a_{\mathrm{obs}}r_u$
		& $0.656904$
		& $<\underline{\Delta}_{\mathrm{obs}}=0.7$\\
		
		\shortstack[l]{
			$\sqrt{N}\,\bar a_{\mathrm c}r_u+\delta_{\mathrm c}$
		}
		& $0.397298$
		& $<\Delta^{\mathrm c}=0.4$\\
		
		\shortstack[l]{
			$\max_{t}\|\tilde u(t)\|$
		}
		& $0.0571577$
		& $<r_u=0.061$\\
		
		\shortstack[l]{
			$\max_{t}
			\delta_{\mathrm c}^{\mathrm{tr}}(t)$
		}
		& $0.00949398$
		& $<\delta_{\mathrm c}=0.03$\\
			\shortstack[l]{
			$\max_{t}
			\|u(t)\|$
		}
		& $13.480099$
		& $<15$\\
		
		\addlinespace[2pt]
		\multicolumn{3}{@{}l}{
			\textit{Theorem \ref{thm:closed_loop_LES}}
		}\\
		\midrule
		
		$\epsilon$
		& $5\times10^{-6}$
		& $<\bar\epsilon=9.31743\times10^{-6}$\\
		
		$\|e_\vartheta(0)\|$
		& $4.45854\times10^{-9}$
		& $<r_{\Theta_1}=5.94472\times10^{-9}$\\
		
		$\|e_f(0)\|$
		& $3.29491\times10^{-7}$
		& $<r_f=6.58982\times10^{-7}$\\
		
		$\eta_{\mathcal I}$
		& $1.53577\times10^{-4}$
		& $>0$\\
		
		\bottomrule
\end{tabularx}
\vspace{2pt}
\parbox{\linewidth}{%
\footnotesize\raggedright
\emph{Note:}
The connectivity minimum is computed along the
nominal HB trajectory, while the tracking-error,
connectivity-discrepancy, and input maxima are
obtained along the tightened trajectory
($t\in[0,1200]$).
Regional residence, QP feasibility, and uniform
error bounds are assumed in Theorem~\ref{thm:regional_safety}.
}
\end{table}
To assess the effect of tightening, the untightened run is
obtained by setting
$
\Delta_{ij}^{\mathrm{int}}
=
\Delta_{i,k}^{\mathrm{obs}}
=
\Delta^{\mathrm c}
=0,
$
while keeping all remaining parameters and the complete
initial state unchanged.
Fig.~\ref{fig:trajectories} compares the resulting trajectories.
The tightened flow avoids collisions and obstacles while
approaching the prescribed target configuration. In contrast,
the untightened flow exhibits obstacle penetration by agent~1
and an inter-agent collision between agents~4 and~5.
\begin{figure}[htbp]
	\centering
	\includegraphics[width=\linewidth]
	{TCNS_final_revision_20261101/trajectory.eps}
	\caption{Agent trajectories under the tightened (solid)
		and untightened (dashed) flows. Circles and diamonds mark initial and target positions.
		Insets A and B show
		the obstacle~$1$ and pairwise~$(4,5)$ constraints, respectively,
		highlighting the violations in the untightened run.}
	\label{fig:trajectories}
\end{figure}



Fig.~\ref{fig:physical-margins} compares the collision
and obstacle margins. The tightened flow maintains
positive safety margins, whereas the untightened flow
violates both collision and obstacle constraints.
Fig.~\ref{fig:connectivity-spectrum} compares the matrix-tree
certificate and algebraic connectivity. Although the
certificate becomes negative in the untightened case,
the algebraic connectivity remains above the prescribed
threshold in both runs. This illustrates that the
matrix-tree condition is sufficient but not necessary
for connectivity preservation.
As shown in Fig.~\ref{fig:edge-switching}, the
communication topology undergoes edge activation and
deactivation, with the number of active edges changing
from $9$ to $10$ and back to $9$.


\begin{figure}[!t]
	\centering
	\subfloat[Collision and obstacle margins.]{
		\includegraphics[width=\linewidth]
		{TCNS_final_revision_20261101/safety.eps}
		\label{fig:physical-margins}
	}
	\\
	\subfloat[Matrix-tree certificate and algebraic connectivity.]{
		\includegraphics[width=\linewidth]
		{TCNS_final_revision_20261101/connectivity.eps}
		\label{fig:connectivity-spectrum}
	}
	\\
	\subfloat[Communication-edge activation and deactivation.]{
		\includegraphics[width=0.95\linewidth]
		{TCNS_final_revision_20261101/edge.eps}
		\label{fig:edge-switching}
	}
	\caption{Safety and connectivity under the
		tightened and untightened flows. Solid and dashed curves
		in (a) and (b) denote the tightened and untightened runs; (c) shows edge switching in the tightened run.}
	\label{fig:safety-margins}
\end{figure}

Finally, we test local convergence using initial conditions
satisfying Theorem~\ref{thm:closed_loop_LES}.
We initialize $u(0)=u_{\mathrm{nom}}(\mathbf y(0))$,
$\|\zeta_0(0)\|=r_f/2$, and all multipliers at zero,
including those associated with the locally omitted
pair $(3,4)$.
The experiment is simulated using \texttt{ode15s}, with
$\mathrm{RelTol}=2\times10^{-10}$,
$\mathrm{AbsTol}=10^{-14}$, and
$\mathrm{MaxStep}=0.1$s.
Fig.~\ref{fig:local-exp} illustrates the decay
of the slow error and the rapid fast transient,
which reaches numerical precision on the
time scale $t/\epsilon$.
The implementation, reference-data links, and
reproduction instructions are available at
\url{https://github.com/SichenQian/safety-critical-gne/tree/v1.0.1}.

\begin{figure}[!h]
	\centering
	\includegraphics[width=\linewidth]
	{TCNS_final_revision_20261101/exp.eps}
\caption{Slow and fast error evolution in an independent
	local experiment satisfying the initial conditions of
	Theorem~\ref{thm:closed_loop_LES}. The inset resolves
	the fast transient on the time scale $t/\epsilon$;
	the subsequent fast-error plateau reflects
	numerical precision.}
	\label{fig:local-exp}
\end{figure}

